% 模型构建模板
% 正文的核心部分——通常是 1-3 个模型的逐一展开
% 每个模型独立为一节

% === 模型应当逐个独立成节 ===
% \section{模型一：<模型名称>}
%   \subsection{模型构建}
%   \subsection{参数估计与求解}
%   \subsection{结果}

\section{模型构建}

% 符号表（如有需要，可放在此处或作为独立符号节）
% 使用 latex/templates/components/notation-table.tex

% <每个模型从对应的内部报告中提取：
%   - 建模思路（为什么这样建）
%   - 数学推导（关键公式，跳过中间步骤）
%   - 参数说明（符号、含义、取值、来源）
%   - 求解方法
%   - 模型结果（数值+图表）>

% 模型之间应有逻辑连接：前面的输出如何成为后面的输入
